\documentclass[12pt,a4paper]{report}

\usepackage{graphicx}
\usepackage{setspace}
\usepackage{geometry}
\usepackage{parskip}
\usepackage{float}
\geometry{left=1in, right=1in, top=1in, bottom=1in}

\begin{document}
\begin{titlepage}
    \centering

    % Institute
    {\Large\bfseries COEP Technological University, Pune}\par
    \vspace{0.4cm}

    % Department (increased font size)
    {\Large Department of Computer Science and Engineering}\par
    \vspace{1.8cm}

    % Title
    {\Huge\bfseries INTERNSHIP REPORT}\par
    \vspace{0.6cm}

    \begin{figure}[H]
        \centering
        \includegraphics[width=0.35\linewidth]{coep_logo.png}
        \label{fig:placeholder}
    \end{figure}

    % Internship info
    \begin{spacing}{1.2}
        \Large
        \begin{tabular}{rl}
            \textbf{Internship Duration:} & 30/05/2025 -- 13/08/2025\\[8pt]
            \textbf{Academic Year:} & TY BTech 2025--26\\
        \end{tabular}
    \end{spacing}

    \vspace{1.6cm}

    % Submitted by
    {\Large\textbf{Submitted by}}\par
    \vspace{0.35cm}
    {\Large \textbf{Chaithanya Gopi Hosahudya}}\par
    \vspace{0.15cm}
    {\large MIS: \textbf{612303041}}\par

    \vspace{1.6cm}

    % Mentors (font size increased)
    \begin{spacing}{1.3}
        \Large
        \textbf{Internal Mentor:} Dr. Suraj T. Sawant\\[10pt]
        \textbf{External Mentor:} Prof. Umakant D. Rapol, IISER Pune
    \end{spacing}

    \vfill
    {\small \today}
\end{titlepage}
\tableofcontents

\chapter*{Introduction}
\addcontentsline{toc}{chapter}{Introduction}

Ion traps are devices that use electric fields to confine charged particles and are widely used in areas such as quantum information processing, high precision spectroscopy and studies involving coherent quantum control. Among the different types of ion traps, radiofrequency (RF) Paul traps are commonly used because they can provide stable confinement through a combination of oscillating and static electric fields. With recent advances in fabrication techniques, surface ion traps have become an important alternative to conventional three dimensional traps. In a surface trap, all electrodes lie on a single plane and the ion is trapped above this surface. These traps are easier to fabricate, more scalable and allow greater flexibility when integrating additional components.

The trap studied in this internship is implemented on a printed circuit board and consists of two RF electrodes placed symmetrically, along with multiple segmented DC electrodes used for axial confinement and fine positioning of the ion. The RF electrodes generate an oscillating quadrupole-like field that produces a pseudopotential minimum above the surface, which provides radial confinement. The DC electrodes shape the potential along the axis of the trap and determine the exact position of the ion.

As part of this work, the fields generated by the RF and DC electrodes were studied in detail using COMSOL Multiphysics. The RF pseudopotential was calculated to identify the height of the trapping region and to examine the nature of radial confinement. The DC fields were analysed to understand the formation of the axial potential well and to study the effect of different voltage settings on the trap behaviour. In addition to these studies, simulations were carried out using the Charged Particle Tracing interface to observe the actual motion of a Ca$^{+}$ ion inside the trap. These simulations provided a clearer picture of the combined effect of the RF and DC fields and helped in checking whether the trap offers stable confinement.

The work also included modelling the RF feed trace on the PCB to understand how the RF signal behaves when delivered to the electrodes. This is important because the performance of the trap depends not only on the electrode geometry but also on how efficiently the RF voltage reaches the trapping region. The information obtained from these simulations forms an initial step towards designing and analysing surface ion traps that can be used in future experimental setups.

\chapter*{Objectives}
\addcontentsline{toc}{chapter}{Objectives}

The main objectives of the internship were as follows:

\begin{itemize}
    \item To study the behaviour of a PCB based surface ion trap through detailed simulations in COMSOL Multiphysics and understand the principles of RF and DC confinement.
    
    \item To calculate the RF pseudopotential, identify the trapping height and examine the radial confinement produced by the RF electrodes.
    
    \item To analyse the DC fields generated by the segmented electrodes and understand the axial potential well and the effect of different voltage settings.
    
    \item To simulate the motion of a Ca$^{+}$ ion inside the trap using the Charged Particle Tracing interface and study the secular and micromotion components.
    
    \item To model the RF feed trace on the PCB and understand how the RF voltage propagates to the electrodes and whether any improvements in the layout may be required.
    
    \item To organise the results, prepare figures and interpret the simulation output in a way that contributes to the ongoing trap design work in the group.
\end{itemize}

\chapter*{About the Organization}
\addcontentsline{toc}{chapter}{About the Organization}

This internship was carried out in the Department of Physics at the Indian Institute of Science Education and Research (IISER) Pune under the guidance of Prof.\ Umakant D.\ Rapol. IISER Pune is a research institute that focuses on science education and advanced research in various fields. The Department of Physics has several research groups that work in areas such as atomic physics, quantum optics, condensed matter and quantum technologies.

Prof.\ Rapol’s group works in experimental atomic physics with a focus on laser cooling, trapped ions, ultracold atoms and precision optical studies. The group is actively working on developing ion trap systems using Ca$^{+}$ ions for quantum information processing and related applications. The work includes trap design, modelling and simulations, vacuum setup, laser systems and electronics for ion control and detection.

IISER Pune also hosts the I-Hub Quantum Technology Foundation (QTF) which supports research and development in quantum technologies. Although this internship was directly under Prof.\ Rapol, the broader research environment at IISER and the goals of QTF have a strong connection with ongoing work in trapped ion systems. This provided a good setting for understanding the practical aspects of ion trap design and the simulation work carried out during this internship.

The simulations and analysis discussed in this report form an initial part of the groundwork for ion trap design and development in the group. The work connects well with the experimental goals of the laboratory and contributes towards the early stages of trap modelling and evaluation.

\chapter*{Scope of Work}
\addcontentsline{toc}{chapter}{Scope of Work}

The main scope of this internship was to study the design and behaviour of a surface ion trap through detailed simulations in COMSOL Multiphysics. The work involved understanding the basic principles of ion trapping and applying these ideas to a specific PCB based trap geometry used in the research group. This included analysing the electric fields, the RF pseudopotential, the DC confinement and the combined effect of these fields on the motion of the ion.

A major part of the work consisted of constructing the three dimensional geometry of the trap in COMSOL and setting up the relevant physics interfaces for the RF, DC and particle tracing studies. This required assigning voltages to the appropriate electrodes, identifying the operating conditions of the trap and calculating the effective potential experienced by the ion. The work also included modelling the RF feed trace on the PCB to obtain an idea of how the RF signal behaves when delivered to the electrodes.

In addition to the field simulations, the scope of the internship included studying the motion of a single Ca$^{+}$ ion using the Charged Particle Tracing interface. This helped in checking whether the trap offers stable confinement and in visualising the secular and micromotion components. Extended structure studies and voltage sweeps were also carried out to understand how additional features on the PCB and different voltage settings could influence the trapping region.

Overall, the scope of work included learning the theoretical aspects of ion traps, becoming familiar with COMSOL as a simulation tool, analysing the results obtained from the simulations and preparing the findings in a clear and structured manner. The work contributes to the early stages of ion trap development in the research group.

\chapter*{Activities Undertaken}
\addcontentsline{toc}{chapter}{Activities Undertaken}

\section*{Understanding the Trap and Theoretical Background}
Before starting the simulations, it was important to understand the basic principles that govern ion trapping. This included reviewing how RF fields produce radial confinement through the pseudopotential approximation and how DC electrodes shape the axial potential. Going through the theory helped in understanding what features to expect in the simulation results and how to interpret the behaviour of the ion in the combined fields. This initial study also clarified the expected trap height, the nature of the RF null and the form of the axial potential.

In addition to the standard treatment, time was spent examining how planar electrode geometries differ from traditional 3D Paul traps. The notes provided in the beginning and earlier sketches of the electrode layout helped in forming a clearer intuition about how the RF null emerges in surface traps and why the confinement is typically weaker along the axial direction. This understanding proved essential when later interpreting the gradual curvature of the DC potential and the characteristic saddle structure around the trapping region.

A review of common microfabricated trap architectures such as the Bastille design on which this PCB trap was loosely modelled also offered helpful reference points. These comparisons guided expectations regarding trap height, field symmetry and the magnitude of micromotion, all of which were reflected closely in the simulation results.

\begin{figure}[H]
  \centering
  \includegraphics[width=0.7\textwidth]{sttv.png}
  \caption{A bastille ion trap.}
\end{figure}

\section*{Geometry Construction}
The first practical step in COMSOL was constructing the geometry of the trap in three dimensions. The trap consisted of two RF rails placed symmetrically and several segmented DC electrodes on a PCB. Care had to be taken to model the electrode widths, gaps and alignment accurately, as even small changes can alter the field near the trapping region. The PCB layer and the surrounding air domain were added to create a realistic model. This geometry served as the foundation for all further simulations.

Beyond the basic layout, certain details required additional consideration. For instance, the copper thickness and the FR4 substrate were both included explicitly to mimic real fabrication conditions. The segmentation of the DC electrodes followed conventions used in several surface traps, helping ensure that the axial potential could be adjusted smoothly without creating unwanted discontinuities. The long routing arms attached to the RF and DC electrodes were also incorporated, since they can introduce small but measurable distortions in the potential if omitted.

Meshing the geometry formed an important part of this stage. The narrow electrode gaps and sharp corners often lead to high electric-field gradients, requiring a refined mesh in these regions. A coarser mesh was used farther from the trap where the fields decay smoothly. This careful combination of local refinement and global efficiency ensured numerical stability in later RF and DC simulations.

\begin{figure}[H]
  \centering
  \includegraphics[width=0.7\textwidth]{c3g.png}
  \caption{Top view of the constructed surface trap geometry.}
\end{figure}

\section*{Simulation Setup}
Once the geometry was ready, the physics interfaces were configured. The RF electrodes were assigned an oscillating voltage for the frequency-domain study, and the DC electrodes were assigned static voltages using the Electrostatics interface. The mesh required refinement in the narrow electrode gaps and along edges where the electric field changes sharply. Several iterations were needed before a stable mesh and solver configuration were obtained. These preparatory steps ensured that the resulting fields were accurate enough to analyse the trap behaviour.

During this phase, a number of diagnostic simulations were run to verify boundary conditions and electrode assignments. Simple electric-field norm plots and potential slices were examined to ensure that the free-space boundary was sufficiently far from the trap and that no artificial distortions arose from the external surfaces. These smaller checks helped avoid unexpected behaviour in later simulations, particularly in the voltage sweeps and particle-tracing studies.

A careful review was also made of the solver configuration, including iterative refinement of the frequency-domain tolerance for the RF study. These adjustments helped stabilise the solution when modelling the strong gradients near the electrode gaps, resulting in smoother pseudopotential distributions.

\section*{RF Pseudopotential Analysis}
The frequency-domain RF simulation gave the oscillating electric fields, from which the pseudopotential was computed. The results showed a distinct minimum above the trap surface, consistent with expectations for a surface RF trap. The trap height and radial confinement could be identified clearly from the contour plots. The symmetry of the pseudopotential also provided an indication that the geometry and boundary conditions were applied correctly.

Additional slices were taken along different planes to examine how the curvature varied in the vicinity of the trap. A line plot along the vertical axis confirmed the expected shallow profile near the minimum, with steep radial rise consistent with planar RF confinement. These analyses also helped identify small asymmetries introduced by the extended copper arms, though such effects remained negligible near the trapping region.

Comparisons with typical pseudopotential shapes from the literature reinforced confidence that the electrode widths and gap dimensions were well chosen. The radial confinement appeared strong enough to support stable trapping when combined with the axial DC confinement studied later.

\begin{figure}[H]
  \centering
  \includegraphics[width=0.7\textwidth]{rfpp_yz.png}
  \caption{RF pseudopotential distribution along the YZ-plane showing the radial confinement region.}
\end{figure}

\section*{DC Potential and Voltage Sweeps}
The DC potential was studied by applying voltages to the segmented electrodes. This produced an axial well whose shape and depth depended on the chosen voltages. Voltage sweeps were carried out to examine how the well shifts or becomes deeper under different configurations. These results are useful when transporting ions along the axis or when creating different trapping zones. The study also confirmed that the segmented layout provides enough flexibility for shaping the axial confinement.

Line cuts of the axial potential helped quantify the trap depth under different voltages. For moderate electrode voltages, the potential well exhibited a smooth and symmetric curvature, providing favourable conditions for stable confinement once the RF contribution was included.

\begin{figure}[H]
  \centering
  \includegraphics[width=0.7\textwidth]{dc_axial.png}
  \caption{Axial DC potential showing the confining well formed by the segmented electrodes.}
\end{figure}

Further analysis showed how individual electrodes particularly the endcap-like EC pads played a dominant role in defining the axial curvature. The middle (ME) and side (SE) electrodes offered finer control, making it possible to adjust the location of the potential minimum without significantly disturbing the radial confinement. These behaviours were visible in potential slices taken just below the electrode plane, where the gradual axial variation could be clearly observed.

\begin{figure}[H]
    \centering
    \begin{minipage}{0.49\linewidth}
        \centering
        \includegraphics[width=\linewidth]{ec.png}
    \end{minipage}
    \begin{minipage}{0.49\linewidth}
        \centering
        \includegraphics[width=\linewidth]{se.png}
    \end{minipage}

    \vspace{0.3cm}

    \begin{minipage}{0.49\linewidth}
        \centering
        \includegraphics[width=\linewidth]{me.png}
    \end{minipage}
    \begin{minipage}{0.49\linewidth}
        \centering
        \includegraphics[width=\linewidth]{ae.png}
    \end{minipage}

    \caption{Representative DC voltage sweep results for the four classes of electrodes: EC, SE, ME and AE.  
    The EC electrodes act as endcap structures in the surface-trap geometry and therefore produce the most pronounced shifts in the axial minimum.  
    The SE and ME electrodes contribute finer shaping of the axial potential, allowing small adjustments to the well curvature and the position of the minimum without disturbing the radial confinement significantly.  
    The AE electrode plays a particularly important role in defining the underlying axial confinement profile; because it spans the central region of the trap, it sets the baseline curvature against which the other segmented electrodes make local modifications.}
    \label{fig:dcsweeps}
\end{figure}


\section*{Combined RF and DC Effective Potential}
The RF pseudopotential and the DC potential were added to obtain the complete effective potential experienced by the ion. This combined potential determines the actual trapping location and the confinement strength in all three dimensions. The resulting plots showed a clear minimum at a stable location above the electrode surface, confirming that the geometry and voltage settings produced a suitable trapping region. This step also provided the background needed to interpret the particle motion seen in later simulations.

The combined view also illustrated how the RF null and the DC saddle do not always coincide perfectly. Small shifts in the axial position of the minimum could be attributed to specific electrode voltages, emphasising the role played by the segmented layout. These variations, though small, informed later choices for particle-tracing initial conditions.

Additional three-dimensional renderings of the potential landscape helped visualise how the trap behaves over an extended region. No stray minima were observed, which is an important requirement for stable operation and for avoiding unintentional ion capture during transport operations.

\begin{figure}[H]
  \centering
  \includegraphics[width=0.7\textwidth]{comb3d.png}
  \caption{Three-dimensional effective potential obtained by combining RF pseudopotential and DC potential.}
\end{figure}

\section*{Charged Particle Tracing}

With the full electric field distribution available from the RF and DC simulations, the next step was to study the actual motion of a single Ca$^{+}$ ion inside the trap. The Charged Particle Tracing interface was used for this purpose, allowing the ion’s trajectory to be computed directly from the simulated fields. The ion was initially placed at the RF null obtained from the pseudopotential analysis, and this starting configuration helped in verifying that the combined fields behaved as expected near the trapping region.

\begin{figure}[H]
  \centering
  \includegraphics[width=0.75\textwidth]{cpt.png}
  \caption{Initial CPT configuration showing the Ca$^{+}$ ion placed at the RF null at $t = 0$.}
\end{figure}

Once the simulation was run, the ion exhibited the characteristic motion associated with RF traps. A slow secular oscillation governed by the effective potential was observed, with a faster micromotion superimposed due to the RF drive. Under stable voltage conditions, the ion remained confined near the RF null throughout the simulation interval, and its motion stayed bounded and centered within the expected trapping region. This behaviour served as an important confirmation that the geometry, meshing and applied voltages were consistent with those required for stable confinement.

Longer simulations were explored to study sensitivity to voltage variations and mesh refinement. When asymmetries were introduced deliberately, the ion trajectory began to drift, eventually gaining enough energy to escape the trapping region. These cases provided useful examples of unstable configurations and highlighted the importance of aligning the RF and DC minima.

\begin{figure}[H]
  \centering
  \includegraphics[width=0.65\textwidth]{traj_xy.png}
  \caption{Simulated Ca$^{+}$ ion trajectory projected in the XY-plane under stable trapping conditions.}
\end{figure}

An animation generated using the CPT results offered an intuitive visualisation of the micromotion-secular motion interplay. Observing the ion's path over many RF cycles also helped confirm that the trapping region was free from higher-order distortions that could cause long-term drifts.

\begin{figure}[H]
  \centering
  \includegraphics[width=0.95\textwidth]{escape.png}
  \caption{Ion escape trajectory showing loss of confinement after approximately $10~\mu\mathrm{s}$.}
\end{figure}

\section*{Extended Structure Analysis}
The PCB used for the trap contains large copper regions beyond the main electrode area. To understand whether these affect the fields near the ion, extended structure simulations were carried out. The results showed only a small influence on the trapping region, indicating that the local electrode geometry mainly determines the confinement. This study gave confidence that practical PCB layouts will not significantly distort the trapping fields.

Additional potential slices were taken near the long copper arms attached to the electrode pads. These slices showed that the potential remains smooth along the arms and that any distortions introduced by them decay rapidly as one moves toward the trap centre. This behaviour confirms that extended routing structures do not compromise trap performance as long as they remain sufficiently symmetric.

\begin{figure}[H]
  \centering
  \includegraphics[width=0.7\textwidth]{ext1.png}
  \caption{Effect of extended copper structures on the trapping region.}
\end{figure}

\section*{Transmission Line Modelling}
The RF signal reaches the trap through a copper trace on the PCB, and its behaviour along this path was studied using a transmission line model. The simulation helped in estimating how much RF amplitude reaches the electrodes and whether any layout modifications might be necessary. Understanding the RF feed behaviour is important for ensuring that the trap operates with the intended confinement strength.

By modelling the copper trace cross-section, parameters such as distributed inductance, resistance and capacitance were extracted. These values indicated that the attenuation along the trace is small, and the RF amplitude at the electrodes should remain close to the intended value. This finding is reassuring for PCB-based traps, where long microstrip lines can sometimes reduce the effective trapping strength if not analysed beforehand.

\begin{figure}[H]
  \centering
  \includegraphics[width=0.7\textwidth]{rftrace.png}
  \caption{Electric field distribution along the RF feed trace on the PCB.}
\end{figure}


\chapter*{Learning and Skill Development}
\addcontentsline{toc}{chapter}{Learning and Skill Development}

This internship provided a good opportunity to learn both the theoretical and practical aspects of ion trap systems. Before starting the simulations, I learned the basic ideas behind RF trapping, the concept of the pseudopotential, the role of DC electrodes and the overall working principle of a linear surface trap. Understanding these ideas helped in setting up the simulations correctly and in interpreting the results in a meaningful way.

A major part of the learning came from working with COMSOL Multiphysics. During the internship, I became familiar with the different physics interfaces used in the project, especially the Electrostatics, Electric Currents and Charged Particle Tracing modules. I learned how to import geometries, assign materials, set boundary conditions and combine the results from different studies. The work also involved creating cut planes, line profiles and potential maps, which helped me understand how to visualise simulation results clearly.

The charged particle tracing study was another important learning experience. Setting up a time dependent particle simulation required choosing a suitable time step, understanding how the fields interact with the particle and ensuring that the solver remains stable throughout the simulation time. This helped me develop a better understanding of how ions actually move inside a trap and how the RF and DC fields influence their motion.

Apart from the technical aspects, the internship also helped me develop general research skills. This included organising simulation data, preparing figures, writing descriptive explanations for the results and understanding how to present the work in a clear and structured way. Overall, the internship provided both theoretical understanding and practical skills that will be useful for future work in this area.

\chapter*{Challenges Faced}
\addcontentsline{toc}{chapter}{Challenges Faced}

During the internship, there were several challenges that I faced while carrying out the simulations and understanding the behaviour of the ion trap. One of the first difficulties was learning how to set up the COMSOL model correctly. The trap geometry contained many narrow gaps and thin copper regions, and meshing these areas without causing solver errors took some time to understand. At the beginning, the model would often fail to solve because of poor mesh quality or incorrect boundary assignments, and this required repeated adjustments until the results became stable.

Understanding the RF simulation was another challenge. Since the RF field is calculated in the frequency domain, interpreting the results and relating them to the time averaged pseudopotential took some effort. It was also difficult at first to combine the RF and DC results correctly and to check whether the values obtained made physical sense. Running the charged particle tracing study was also not straightforward. The particle would sometimes escape the trap very quickly, and I had to review the solver settings, time step and initial conditions before the trajectory became stable.

Working with large geometries also caused memory issues in COMSOL. Some of the extended structure simulations would take a long time to run or fail due to insufficient memory. This required reducing the mesh size or simplifying parts of the model so that the simulation could complete. Although this slowed down the work at times, it helped me understand how to manage resources more effectively.

Finally, interpreting the results and connecting them with the underlying physics also took some time. The connection between the RF null, the DC saddle and the trajectory of the ion was not immediately obvious, and understanding these relationships required going back to the theory and studying the plots carefully.

In the end, these challenges played an important role in helping me understand the simulations better and in improving the way I approached the work. Each difficulty required me to slow down, revisit the theory and think more carefully about the steps involved. This gradually made me more comfortable with the tools and the underlying ideas, and the experience has added to my overall learning during the internship.

\chapter*{Impact Assessment}
\addcontentsline{toc}{chapter}{Impact Assessment}

The work carried out during this internship contributed to the initial stages of ion trap design in the group. The simulations helped in understanding how the proposed PCB based trap behaves under different voltage conditions and whether the geometry is suitable for stable ion confinement. The results from the RF and DC studies provided a clear picture of the potential landscape, and the combined analysis helped in confirming the approximate trap height and the overall structure of the effective potential.

The charged particle tracing simulations added an important element to the study. By observing the actual motion of the ion in the simulated fields, it became possible to check the stability of the trap in a more direct way. These results can be useful when comparing the simulated behaviour with future experimental observations. The animation created from the simulation data also offers a simple and clear way to visualise how the ion moves in the trap.

The transmission line study gave additional information about how the RF signal would behave on the PCB. This can help in understanding how much RF voltage will reach the electrodes and whether any changes in the routing or trace width might be required in future designs. The extended structure analysis also helped in checking whether the large copper regions on the board introduce any unwanted effects near the trapping region.

Overall, the work provided useful information that can support the next steps in the development of the trap. The results can be taken forward when refining the geometry, planning the fabrication or preparing the experimental setup. The study also helps in building a clearer understanding of how surface traps behave and what factors must be considered when designing them.

\chapter*{Conclusion}
\addcontentsline{toc}{chapter}{Conclusion}

The simulations carried out during this internship provided a comprehensive understanding of the behaviour of a PCB based surface ion trap, beginning from the basic electrode geometry and extending to the dynamical motion of a single trapped ion. By modelling the geometry in COMSOL and analysing both the RF and DC contributions independently, it was possible to construct a clear picture of the effective trapping potential and to verify that the design is capable of producing stable confinement at a height that matches expectations for surface traps of this size. The RF pseudopotential analysis confirmed that the trap forms a well-defined minimum above the surface, with strong radial curvature arising from the RF electrodes.

The DC potential studies demonstrated that the segmented electrode layout offers sufficient flexibility to shape the axial potential and to shift the trapping zone smoothly along the axis by adjusting the voltages applied to individual electrodes. The combined RF and DC potentials produced a three dimensional well that is smooth and free from irregularities, and the electrode sweeps helped in identifying configurations that strengthen or weaken the confinement. These results suggest that precise control over ion position and stability is achievable in this design.

The ion trajectory simulations provided a more detailed view of the trap’s behaviour. The Ca$^{+}$ ion exhibited both micromotion and secular motion characteristic of RF traps, and in most cases the trajectories remained bounded for the entire simulation. The animation created from these trajectories offered a visual confirmation of the ion’s stability and showed how the combined potential governs its motion. Although a few voltage settings resulted in ion loss, these cases helped in identifying regions of operation that may be avoided in practice.

Additional studies on the extended copper structures and the RF feed trace showed that the trap behaves well even when its surrounding geometry is taken into account. The transmission line modelling indicated that the RF signal should reach the electrodes with minimal attenuation, and the extended structure analysis confirmed that distant copper regions do not introduce significant distortions near the trapping region. These findings are important for practical implementation, as they support the suitability of the PCB layout for fabrication.

Overall, the work completed during this internship forms a useful starting point for further development of the trap. The results provide guidance for refining the design, planning the fabrication and preparing the experimental setup. The study also contributes to a better understanding of surface trap behaviour, which will be valuable for future work in ion trap based quantum technologies.

\chapter*{Recommendations}
\addcontentsline{toc}{chapter}{Recommendations}

Based on the work carried out during this internship, several points may be considered for future studies. One area that can be explored further is the behaviour of the trap under a wider range of DC voltage settings. The present analysis focused mainly on single electrode sweeps, and extending this to combinations of electrodes may give a clearer understanding of how the axial potential can be shaped for operations such as transport or splitting of ions. This may also help in checking how sensitive the trap is to small changes in voltage.

It may also be useful to refine the geometry of the trap in the simulation by including additional details that were not modelled in the present work. Features such as stitching vias, variations in copper thickness or small asymmetries that may arise during fabrication can sometimes influence the potential near the trapping region. Including these in the simulation would provide a more accurate picture of the actual PCB and help in identifying any corrections that may be needed before fabrication.

The transmission line model can also be extended. In the current study, a simple straight trace was considered. Studying the effect of bends, different trace widths or the presence of nearby ground planes may help in understanding how the RF voltage behaves on more complex layouts. This could be important when the trap is used in larger systems or when additional components are placed on the same board.

Further improvements may be made to the charged particle tracing studies. Running longer simulations or using different starting positions and energies may help in understanding long-term stability and the effect of small perturbations. It may also be helpful to simulate the trajectories under slightly different RF amplitudes or frequencies to see how the trap behaves under non-ideal conditions. These studies can provide useful information when comparing simulation results with the behaviour observed in the actual experiment.

Finally, as the design develops, it may be beneficial to compare the simulation results with measurements from a fabricated version of the trap. This would help in checking how closely the simulation matches reality and whether any adjustments are needed in the modelling approach.

\chapter*{Acknowledgements}
\addcontentsline{toc}{chapter}{Acknowledgements}

I would like to express my sincere thanks to Prof.\ Umakant D.\ Rapol for giving me the opportunity to work on this project and for his guidance throughout the internship. His explanations and suggestions helped me understand the subject more clearly and provided useful direction at various stages of the work.

I am grateful to Dr.\ Jay Mangaonkar for taking the time to discuss the fundamentals of trapped ion systems with me, for showing me around the lab and for guiding me during the initial phase of the internship. His discussions helped me build the foundation required for the rest of the work.

I would also like to thank Ms.\ Shuvarati Roy for her help while I was learning to use COMSOL. Her experience and inputs made it easier for me to understand the software and set up the initial models correctly.

I am especially thankful to Mr.\ Rayees A.\ S.\ for his continuous support throughout the internship. His willingness to spend time discussing the simulations, helping me troubleshoot problems and clarifying concepts played an important role in completing the work. I am sincerely appreciative of the guidance he provided during this period.

I would also like to thank the members of the research group for their support and for the discussions that helped me look at the work from different perspectives. I am thankful to IISER Pune for providing the facilities required to carry out the simulations.

I would also like to express my gratitude to my internal mentor, Dr. Suraj T.\ Sawant, for his guidance and support throughout the duration of the internship. I am equally thankful to COEP Technological University for permitting me to undertake this internship as part of the curriculum and for providing the opportunity to gain meaningful experience in a research environment.

Lastly, I would like to thank my family for their support during the course of this internship.

\chapter*{References}
\addcontentsline{toc}{chapter}{References}

\begin{enumerate}
    \item COMSOL Multiphysics Documentation, AC/DC Module and Charged Particle Tracing Module.
    \item IISER Pune, Department of Physics, information available on the institute website.
    \item Atomic Physics and Quantum Optics Lab (Prof.\ Umakant D.\ Rapol), IISER Pune, Research overview and materials.
    \item Standard material on RF ion traps, pseudopotential approximation and trapping principles.
    \item Figures, data and simulation results generated during the internship work.
\end{enumerate}


\end{document}
